\begin{notechapter}
\chapter{Bidirectional RNNs and Autoregressive Sequence Models}\label{ch:sequence-models}
\chaptersubtitle{What a recurrent state summarizes, and how to choose the direction}

\begin{chapterintro}
The state of a forward RNN at time $t$ is a function of $x_{1:t}$: it summarizes the past. Many
problems need something else. To tag a word, the words that follow it matter as much as
those before it. To infer the hidden state of a dynamical system from a whole recording, the
future observations matter. To define a probability distribution over sequences, the
model must predict each element from the previous ones only, and never see the element it
predicts. This chapter shows how the same recurrent cell, run in different directions on
differently shifted inputs, produces each of these summaries. You will learn
\begin{itemize}[nosep]
  \item that a forward state summarizes the prefix $x_{1:t}$ and a backward state the suffix $x_{t:T}$;
  \item how a backward RNN is implemented by flipping the sequence, and why this needs care
        with padded batches;
  \item bidirectional RNNs and the ways of combining their two directions;
  \item autoregressive models: how an RNN on the \emph{shifted} sequence defines a normalized
        probability distribution over sequences, how it is trained and how it generates;
  \item why a model that sees the element it predicts is not a probability model.
\end{itemize}
\prerequisites{\cref{ch:rnn-bptt,ch:gated-rnn}; the chain rule of probability and maximum
likelihood (\cref{ch:cross-entropy}).}
\end{chapterintro}

\section{What a recurrent state summarizes}\label{sequence-models:sec:summaries}

Write the recurrence of any of the previous chapters abstractly as $h_t = F(h_{t-1}, x_t)$, with
$h_0$ fixed and $F$ the cell (simple, LSTM or GRU; for the LSTM the state is the pair $(h_t, c_t)$).

\begin{proposition}[Forward and backward summaries]\label{sequence-models:prop:summaries}
\begin{enumerate}[nosep]
  \item The forward state $h_t = F(h_{t-1}, x_t)$, $h_0$ fixed, is a function of $x_{1:t}$ only.
  \item The backward state $b_t = F(b_{t+1}, x_t)$, computed for $t = T,\dots,1$ from a fixed $b_{T+1}$, is a function of
        $x_{t:T}$ only.
\end{enumerate}
\end{proposition}

\begin{proof}
(1) By induction: $h_0$ is a constant, and if $h_{t-1}$ is a function of $x_{1:t-1}$, then $F(h_{t-1}, x_t)$ is a function of
$x_{1:t}$. (2) The same induction, downward from $b_{T+1}$.
\end{proof}

The proposition says what each state \emph{can} depend on. Which information it actually
keeps from its prefix or suffix is learned, within the limits discussed in
\cref{rnn-bptt:sec:vanishing}. At $t = T$ the backward state has seen only $x_T$; at $t = 1$ it has seen the
whole sequence. The forward state is the opposite.

\section{The backward RNN by flipping}\label{sequence-models:sec:flip}

Libraries implement recurrences that run forward in time. A backward RNN is obtained from
a forward one in three steps:
\begin{enumerate}[nosep]
  \item flip the sequence along the time axis: $(x_T, x_{T-1},\dots,x_1)$;
  \item run a forward RNN on it: its $k$-th output $h'_k$, $k = 1,\dots,T$, has seen the first $k$ elements of the flipped
        sequence, $x_T,\dots,x_{T-k+1}$;
  \item flip the outputs back: the output now stored at position $t = T - k + 1$ has seen $x_{T-k+1},\dots,x_T$, that is $x_{t:T}$.
\end{enumerate}
The output at position $t$ after the second flip is therefore $b_t$, and the recursion is the one of
\cref{sequence-models:prop:summaries} (2), since $h'_k = F(h'_{k-1}, x_{T-k+1})$ reads $b_t = F(b_{t+1}, x_t)$ with $t = T - k + 1$.

\begin{figure}[ht]
\centering
\begin{tikzpicture}[
  cell/.style={draw, minimum width=9mm, minimum height=7mm, inner sep=1pt, font=\small},
  ar/.style={-{Stealth[length=1.8mm]}, thick},
  lab/.style={font=\scriptsize, inner sep=1pt}]
  \def\cA{bookaccent!30}\def\cB{bookexample!35}\def\cC{bookdef!35}\def\cD{bookpitfall!25}
  % row 1
  \node[lab, anchor=east] at (-0.2,0) {sequence};
  \foreach \i/\c in {1/\cA,2/\cB,3/\cC,4/\cD} \node[cell, fill=\c] (a\i) at (0.4+1.4*\i,0) {$x_{\i}$};
  % row 2
  \node[lab, anchor=east] at (-0.2,-1.5) {flipped};
  \foreach \k/\i/\c in {1/4/\cD,2/3/\cC,3/2/\cB,4/1/\cA} \node[cell, fill=\c] (b\k) at (0.4+1.4*\k,-1.5) {$x_{\i}$};
  \draw[ar, gray] (3.9,-0.45) -- node[lab, right]{\texttt{torch.flip(x, [0])}} (3.9,-1.05);
  % row 3: forward RNN
  \node[lab, anchor=east] at (-0.2,-3.0) {forward RNN};
  \foreach \k in {1,2,3,4} {
    \node[cell, fill=gray!10] (c\k) at (0.4+1.4*\k,-3.0) {$h'_{\k}$};
    \draw[ar] (b\k) -- (c\k);
  }
  \foreach \k/\j in {1/2,2/3,3/4} \draw[ar, bookaccent] (c\k) -- (c\j);
  \node[lab, below=1pt of c1] {$x_{4}$};
  \node[lab, below=1pt of c2] {$x_{3:4}$};
  \node[lab, below=1pt of c3] {$x_{2:4}$};
  \node[lab, below=1pt of c4] {$x_{1:4}$};
  % row 4: flipped back
  \node[lab, anchor=east] at (-0.2,-5.0) {flipped back};
  \foreach \t/\k in {1/4,2/3,3/2,4/1} \node[cell, fill=bookkey!12] (d\t) at (0.4+1.4*\t,-5.0) {$b_{\t}$};
  \draw[ar, gray] (3.9,-3.75) -- node[lab, right]{\texttt{torch.flip(h, [0])}} (3.9,-4.55);
  \node[lab, below=1pt of d1] {$x_{1:4}$};
  \node[lab, below=1pt of d2] {$x_{2:4}$};
  \node[lab, below=1pt of d3] {$x_{3:4}$};
  \node[lab, below=1pt of d4] {$x_{4}$};
  % right annotations
  \node[lab, align=left, anchor=west] at (6.6,-3.0) {$h'_k$ has seen the first $k$\\ elements of the flipped\\ sequence};
  \node[lab, align=left, anchor=west] at (6.6,-5.0) {$b_t = h'_{T-t+1}$ summarizes\\ the suffix $x_{t:T}$};
\end{tikzpicture}
\caption{A backward RNN implemented with a forward one, for $T = 4$. Below each output: the inputs it
has seen. After the second flip, position $t$ holds a summary of $x_{t:T}$.}
\label{sequence-models:fig:flip}
\end{figure}

This is the module used by the inference network of the deep Kalman filter
(\cref{ch:deep-kalman-filters}), from the repository \texttt{CL-BZH/Deep-Kalman-Filters}; tensors have shape
$(T, \text{batch}, \text{features})$, so time is dimension $0$:
\begin{lstlisting}[language=Python]
class DKF_Backward_RNN(nn.Module):
    def __init__(self, input_size, hidden_size, num_layers):
        super().__init__()
        self.bw_rnn = nn.LSTM(input_size, hidden_size, num_layers)

    def forward(self, x):                      # x: (seq_len, batch, input_size)
        x = torch.flip(x, [0])                 # (x_T, ..., x_1)
        bw_rnn_out, _ = self.bw_rnn(x)         # position k: x_T..x_{T-k+1}
        return torch.flip(bw_rnn_out, [0])     # position t: x_t..x_T
\end{lstlisting}
Two checks confirm the bookkeeping. First, PyTorch's \texttt{nn.LSTM(\ldots, bidirectional=True)} has
a reverse direction; copying its reverse weights into the forward LSTM above gives exactly
the second half of its output. Second, perturbing $x_1$ changes $b_1$ and leaves $b_2,\dots,b_T$ unchanged,
as \cref{sequence-models:prop:summaries} requires.

\begin{pitfall}[Flipping a padded batch]
Sequences of different lengths are batched by padding them to a common length. Flipping a
padded tensor puts the padding \emph{first}: a sequence of length $3$ padded to $5$ becomes
$(0, 0, x_3, x_2, x_1)$, and its backward states are computed after reading two padding vectors, so
they differ from those of the unpadded sequence. PyTorch's bidirectional modules handle this
correctly when the batch is given as a \texttt{PackedSequence} (\texttt{pack\_padded\_sequence}), in
which case the reverse direction starts at each sequence's own last element; on a padded
tensor they do not. The DKF uses sequences of equal length, so the problem does not arise
there.
\end{pitfall}

\section{Bidirectional RNNs}\label{sequence-models:sec:bidirectional}

A \emph{bidirectional} RNN (Schuster and Paliwal, \emph{Bidirectional Recurrent
Neural Networks}, 1997) runs a forward and a backward RNN, with their
own weights, over the same sequence, and combines their states at each position:
\[
  \overrightarrow h_t = F_{\to}(\overrightarrow h_{t-1}, x_t),\qquad
  \overleftarrow h_t = F_{\leftarrow}(\overleftarrow h_{t+1}, x_t),\qquad
  s_t = \mathrm{combine}(\overrightarrow h_t, \overleftarrow h_t) .
\]
By \cref{sequence-models:prop:summaries}, $s_t$ is a function of the whole sequence $x_{1:T}$, with an emphasis on the
neighbourhood of $t$. Common combinations are the concatenation $s_t = (\overrightarrow h_t, \overleftarrow h_t)\in\R^{2m}$ (the PyTorch
default), the sum, a weighted sum $\alpha\overrightarrow h_t + (1 - \alpha)\overleftarrow h_t$, and a linear map of the concatenation;
the \texttt{DKF\_Bidirectional\_RNN} module of the DKF repository offers these four options.

\begin{figure}[ht]
\centering
\begin{tikzpicture}[
  st/.style={draw, circle, minimum size=7mm, inner sep=0pt, font=\scriptsize},
  io/.style={draw, rectangle, rounded corners=2pt, minimum size=6.5mm, inner sep=1pt, font=\scriptsize},
  ar/.style={-{Stealth[length=1.8mm]}, thick}]
  \foreach \t in {1,2,3,4} {
    \node[io, minimum width=11mm] (x\t) at (2*\t+0.2,0) {$x_{\t}$};
    \node[st, fill=bookaccent!15] (f\t) at (2*\t,1.3) {$\overrightarrow{h}_{\t}$};
    \node[st, fill=bookexample!20] (b\t) at (2*\t+0.7,2.5) {$\overleftarrow{h}_{\t}$};
    \node[io, fill=bookkey!12] (s\t) at (2*\t+0.35,3.8) {$s_{\t}$};
    \draw[ar] (x\t) -- (f\t);
    \draw[ar] (2*\t+0.7,0.33) -- (b\t);
    \draw[ar, gray] (f\t) -- (s\t);
    \draw[ar, gray] (b\t) -- (s\t);
  }
  \foreach \t/\u in {1/2,2/3,3/4} {
    \draw[ar, bookaccent] (f\t) -- (f\u);
    \draw[ar, bookexample!80!black] (b\u) -- (b\t);
  }
\end{tikzpicture}
\caption{A bidirectional RNN. The forward chain (blue) runs left to right, the backward chain
(green) right to left, with separate weights; $s_t$ combines $\overrightarrow h_t$ (a summary of $x_{1:t}$) and $\overleftarrow h_t$ (a summary
of $x_{t:T}$).}
\label{sequence-models:fig:birnn}
\end{figure}

A bidirectional RNN is the right tool whenever the whole sequence is available before any
output is needed: tagging the words of a sentence (\cref{ch:hmm-viterbi,ch:crf} treat the
same task with structured models), speech recognition on recorded audio, or inferring the
hidden states of a recorded trajectory, as in \cref{ch:deep-kalman-filters}. It cannot be used to
predict the future, nor, as the next section shows, to define a probability distribution
over sequences.

\section{Autoregressive sequence models}\label{sequence-models:sec:autoregressive}

\paragraph{The chain rule of probability.} Any joint distribution of a sequence factorizes as
\[
  p(x_1,\dots,x_T) = p(x_1)\,p(x_2\mid x_1)\,p(x_3\mid x_1, x_2)\cdots p(x_T\mid x_{1:T-1}) = \prod_{t=1}^Tp(x_t\mid x_{<t}),
\]
with $x_{<t} = x_{1:t-1}$ ($x_{<1}$ is empty). A model of the sequence is therefore a model of each conditional: a
\emph{density that depends on the past}. The difficulty is that the conditioning set grows
with $t$. An RNN solves it: a forward state summarizes the past in a vector of fixed size.

\paragraph{The shifted input.} To predict $x_t$, the state must have seen $x_{<t}$ and \emph{not} $x_t$. The
forward state $h_t$ has seen $x_{1:t}$, one element too many. The model therefore runs the RNN on the
sequence shifted by one step, $(x_0, x_1,\dots,x_{T-1})$, where $x_0$ is a fixed start symbol (a zero vector,
a learned vector, or a special token in text):
\[
  h_0 = F(h_{\mathrm{init}}, x_0),\qquad h_t = F(h_{t-1}, x_t)\ (t\geq1),\qquad
  p_\theta(x_t\mid x_{<t}) = \pi\bigl(x_t;\,G(h_{t-1})\bigr),
\]
where $G$ is an output layer and $\pi(\cdot;\,\eta)$ a distribution with parameters $\eta$: a softmax over a vocabulary for
discrete symbols, a Gaussian with mean and log-variance $G(h_{t-1})$ for real vectors. By
\cref{sequence-models:prop:summaries}, $h_{t-1}$ is a function of $x_{0:t-1}$, so the conditional of $x_t$ depends on the past only.

\begin{proposition}[An autoregressive RNN is a normalized model]\label{sequence-models:prop:normalized}
For discrete $x_t$ and any parameters $\theta$, $\sum_{x_{1:T}}\prod_{t=1}^Tp_\theta(x_t\mid x_{<t}) = 1$. For densities, the same holds with
integrals.
\end{proposition}

\begin{proof}
Sum over $x_T$ first: only the last factor depends on $x_T$, and it sums to $1$ because it is a
distribution over $x_T$ for each fixed $x_{<T}$. What remains is the same expression with $T - 1$ in place
of $T$; by induction the total is $1$. For densities the sums are integrals of nonnegative
functions, which can be taken in this order (Tonelli's theorem).
\end{proof}

The proof uses one property of each factor: it is a distribution over $x_t$ \emph{that does not depend on
$x_t$}. The shifted input is what guarantees it.

\begin{figure}[ht]
\centering
\begin{tikzpicture}[
  st/.style={draw, circle, minimum size=7mm, inner sep=0pt, font=\scriptsize, fill=bookaccent!15},
  io/.style={draw, rectangle, rounded corners=2pt, minimum height=6.5mm, minimum width=8mm, inner sep=1pt, font=\scriptsize},
  ar/.style={-{Stealth[length=1.8mm]}, thick}]
  \foreach \t/\prev in {1/0,2/1,3/2,4/3} {
    \node[io] (x\t) at (2.7*\t,0) {$x_{\prev}$};
    \node[st] (h\t) at (2.7*\t,1.3) {$h_{\prev}$};
    \node[io, fill=bookkey!12] (p\t) at (2.7*\t,2.7) {$p_\theta(x_{\t}\mid x_{<\t})$};
    \node[io, fill=bookexample!15] (y\t) at (2.7*\t,4.0) {$x_{\t}$};
    \draw[ar] (x\t) -- (h\t); \draw[ar] (h\t) -- node[right, font=\scriptsize]{$G$} (p\t);
    \draw[ar, gray, dashed] (y\t) -- node[right, font=\scriptsize]{target} (p\t);
  }
  \foreach \t/\u in {1/2,2/3,3/4} \draw[ar, bookaccent] (h\t) -- (h\u);
  \foreach \t/\u in {1/2,2/3,3/4} \draw[ar, bookpitfall, densely dotted] (y\t.east) -- ++(0.95,0) |- (x\u.west);
\end{tikzpicture}
\caption{An autoregressive RNN for $T = 4$. The input is the sequence shifted by one step, with a start
symbol $x_0$, so $h_{t-1}$ has seen $x_{0:t-1}$ and the output layer $G$ gives the conditional distribution of $x_t$. In
training (teacher forcing) the inputs are the observed $x_t$; in generation (dotted) each sampled $x_t$
becomes the next input.}
\label{sequence-models:fig:autoregressive}
\end{figure}

\paragraph{Training.} The negative log-likelihood of an observed sequence is
\[
  -\log p_\theta(x_{1:T}) = -\sum_{t=1}^T\log p_\theta(x_t\mid x_{<t}),
\]
a sum of per-step losses: cross-entropies for a softmax output, Gaussian negative
log-likelihoods for a Gaussian output (which reduce to squared errors, up to constants, when
the variance is fixed). It is minimized by BPTT (\cref{ch:rnn-bptt}). During training the inputs are
the \emph{observed} previous elements, all available at once, so the whole sequence is
processed in one pass; this is called \emph{teacher forcing} (Williams and Zipser, \emph{A Learning
Algorithm for Continually Running Fully Recurrent Neural Networks}, 1989).

\paragraph{Generation and forecasting.} To sample a sequence, draw $x_1$ from $p_\theta(\cdot\mid x_0)$, feed it as
the next input, draw $x_2$, and so on: ancestral sampling, one element at a time. To forecast
the continuation of an observed sequence, run the RNN on the observed part, then continue
in the same way. At generation time the model is conditioned on its own samples, which can
drift away from the data it was trained on; the mismatch between training on observed
prefixes and generating from sampled ones is called \emph{exposure bias}, and
\emph{scheduled sampling} reduces it by progressively replacing observed inputs with the
model's own predictions during training (Bengio et al., \emph{Scheduled Sampling for Sequence
Prediction with Recurrent Neural Networks}, 2015).

\begin{example}[A binary autoregressive RNN]\label{sequence-models:ex:binary}
Let $x_t\in\{0, 1\}$, $T = 3$, start symbol $x_0 = 0$, a scalar state $h_t = \tanh(0.5\,h_{t-1} + 2x_t)$ with $h_{-1} = 0$ (so $h_0 = 0$), and
$p(x_t = 1\mid x_{<t}) = \sigma(3h_{t-1} - 1)$. For $x = (1, 0, 1)$:
\begin{center}
\begin{tabular}{c|ccc}
$t$ & $h_{t-1}$ & $p(x_t = 1\mid x_{<t})$ & factor\\\hline
$1$ & $0$ & $\sigma(-1) = 0.2689$ & $0.2689$\\
$2$ & $\tanh(2) = 0.9640$ & $\sigma(1.8921) = 0.8690$ & $1 - 0.8690 = 0.1310$\\
$3$ & $\tanh(0.4820) = 0.4479$ & $\sigma(0.3438) = 0.5851$ & $0.5851$
\end{tabular}
\end{center}
so $p(1, 0, 1) = 0.2689\times0.1310\times0.5851 = 0.0206$ and $-\log p = 3.8818$. The eight sequences have probabilities
\[
  \begin{array}{cccccccc}
  000 & 001 & 010 & 011 & 100 & 101 & 110 & 111\\
  0.3907 & 0.1437 & 0.0258 & 0.1709 & 0.0146 & 0.0206 & 0.0289 & 0.2048
  \end{array}
\]
which sum to $1$, as \cref{sequence-models:prop:normalized} requires. The model has learned (or rather, its
weights encode) that a $1$ makes another $1$ likely: $p(111)$ is ten times $p(101)$.
\end{example}

\begin{pitfall}[A model that sees what it predicts]
Run the same RNN on the \emph{unshifted} sequence and use $\sigma(3h_t - 1)$, where $h_t$ has seen $x_t$, as
``$p(x_t = 1\mid\cdot)$''. The product of these factors is no longer a probability distribution over
sequences: in the example above the eight products sum to $3.18$. Such a model can make
its training loss arbitrarily small by learning to copy its input, and it cannot generate,
since it needs $x_t$ to predict $x_t$. The same holds for a bidirectional RNN, whose state at $t$ has seen
the whole sequence. A sequence model is a probability model only if each conditional
ignores the element it predicts and everything after it.
\end{pitfall}

\section{Which summary for which conditional}\label{sequence-models:sec:which}

The same recurrent cell gives four different summaries, depending on the direction and
the shift:
\begin{center}\small
\begin{tabular}{p{0.3\textwidth}p{0.27\textwidth}p{0.33\textwidth}}
\textbf{RNN} & \textbf{state at $t$ summarizes} & \textbf{typical use}\\\hline
forward, on $(x_0,\dots,x_{T-1})$ & $x_{0:t-1}$, the strict past & $p(x_t\mid x_{<t})$: generation, forecasting\\
forward & $x_{1:t}$, the past and present & filtering, online tagging\\
backward & $x_{t:T}$, the present and future & smoothing (with a past summary)\\
bidirectional & $x_{1:T}$, everything & offline tagging and inference
\end{tabular}
\end{center}
Choosing the summary is choosing what each conditional distribution may depend on, and
that choice should follow the conditional independences of the model. The deep Kalman filter
of the next chapter is a complete example: its generative model needs a summary of the
strict past, and its inference model a summary of the future.

\begin{chaptersummary}
\begin{itemize}
  \item A forward recurrent state is a function of $x_{1:t}$, a backward one of $x_{t:T}$
        (\cref{sequence-models:prop:summaries}).
  \item A backward RNN is a forward RNN applied to the flipped sequence, with its outputs flipped
        back; with padded batches, use packed sequences.
  \item A bidirectional RNN combines both summaries; it suits offline tasks where the whole
        sequence is known.
  \item An autoregressive RNN runs on the sequence shifted by one step, so that $h_{t-1}$ summarizes $x_{<t}$
        and parameterizes $p_\theta(x_t\mid x_{<t})$; the product of these conditionals is normalized for any
        parameters. It is trained with teacher forcing and generates by feeding its samples
        back.
  \item A conditional that sees the element it predicts does not define a distribution over
        sequences.
\end{itemize}
\end{chaptersummary}

\section*{Exercises}
\addcontentsline{toc}{section}{Exercises}

\begin{exercise}\label{sequence-models:exr:1}
For $T = 5$, list for each position $k$ of the flipped sequence the inputs seen by $h'_k$, and the
position $t$ at which $h'_k$ is stored after the second flip.
\solutionlink{sequence-models:exr:1}
\end{exercise}

\begin{exercise}\label{sequence-models:exr:2}
For the LSTM, the recurrent state is the pair $(h_t, c_t)$. Check that \cref{sequence-models:prop:summaries} still
applies, and conclude that the output $b_t$ of \texttt{DKF\_Backward\_RNN} depends on $x_{t:T}$ only.
\solutionlink{sequence-models:exr:2}
\end{exercise}

\begin{exercise}\label{sequence-models:exr:3}
In \cref{sequence-models:ex:binary}, compute $p(1, 1, 1)$ and check the value $0.2048$.
\solutionlink{sequence-models:exr:3}
\end{exercise}

\begin{exercise}\label{sequence-models:exr:4}
Take $T = 1$ and the ``leaky'' model $q(x_1 = 1\mid x_1) = \sigma(3\tanh(2x_1) - 1)$. Compute $q(0\mid 0) + q(1\mid 1)$. Why is this
not $1$?
\solutionlink{sequence-models:exr:4}
\end{exercise}

\begin{exercise}\label{sequence-models:exr:5}
A Gaussian autoregressive RNN outputs $p_\theta(x_t\mid x_{<t}) = \mathcal N(\mu_t, \sigma^2)$ with $\mu_t = G(h_{t-1})$ and a fixed $\sigma$.
Show that maximizing the likelihood is equivalent to minimizing $\sum_t\norm{x_t - \mu_t}^2$. What changes if the
network also outputs $\log\sigma_t^2$?
\solutionlink{sequence-models:exr:5}
\end{exercise}
\end{notechapter}
